% !TEX root = ../main.tex
% ============================================================
% ANEXO A. Tarea 86 del cronograma, pedida por el autor el 21/10:
% contrasta el párrafo "Dónde rinde más el MPASM" de la discusión.
% ============================================================
\chapter{Propagación a distancias largas: interferometría holográfica simulada}
\label{anx:larga}
% Repo: scripts/interferometria_larga.py -> resultados/interferometria_larga/
%       resumen.csv (21/10, CPU, complex128). Figura:
%       scripts/figuras_manuscrito.py (figura_interferometria).

La discusión (Sección~\ref{sec:discusion}) sostiene que el MPASM resulta más
útil en montajes de propagación larga, como la interferometría holográfica
digital, porque su $K_f$ automático crece con la distancia. Este anexo lo
contrasta en una geometría de ese tipo, simulada, y separa, como el resto del
trabajo, la propagación de la reconstrucción.

\paragraph{El ensayo.} Un sensor de $1024\times1024$ píxeles de
\SI{3.45}{\micro\metre} registra, a \SI{532}{nm}, la luz de un objeto que
ocupa una ventana del mismo tamaño, coaxial con él. El objeto es un difusor:
un campo complejo aleatorio con el espectro limitado a ángulos de
\SI{10}{mrad} ---manchas de \emph{speckle} de unos \SI{53}{\micro\metre}---,
sobre un disco de \SI{3.2}{mm} de diámetro con borde suave. Como en la
interferometría holográfica, se registra en dos estados, y el segundo lleva la
deformación de una placa circular empotrada: una diferencia de fase
$\Delta\varphi(r) = \Delta\varphi_\mathrm{max}(1 - r^2/a^2)^2$ dentro del
disco de radio $a$, con $\Delta\varphi_\mathrm{max} = 3\pi$. La distancia se
da por el cociente de la ventana, $q = z/(N\delta^2/\lambda)$, entre 0.5 y 20:
con este sensor, de \SI{11}{mm} a \SI{0.46}{m}; con uno de 2048 píxeles, el
mismo $q$ cae al doble de distancia, hasta \SI{0.92}{m}.

La referencia es el espectro angular en una malla rellenada lo bastante para
que la luz que sale de la ventana no reentre y la función de transferencia
quede bien muestreada en toda la banda del difusor. Repetida con 1.3 veces ese
relleno, cambia menos de $6\times10^{-9}$ a la ida y de $2\times10^{-6}$ a la
vuelta. Antes de devolverlo al objeto, el holograma se apodiza con una ventana
de Tukey, la misma para todos los métodos: sin ella, el borde del sensor
difracta hasta la frecuencia de Nyquist y, a $q = 20$, esa luz se abre más de
lo que contiene la propia referencia. Se comparan el FFT-ASM sin relleno, con
relleno al doble y al cuádruple; el BLAS con relleno al doble; el MPASM con
$s = 1$, 2 y 4 y el $K_f$ automático, y, a la vuelta, la transformada de
Fresnel de una sola FFT, que es paraxial y tiene malla de salida
$\lambda z/(N\delta)$. Se mide el error relativo del campo en el sensor
(Ec.~\eqref{eq:error-relativo}), que es el de simular el holograma, y el RMS de
la diferencia de fase reconstruida frente a la deformación, pesado por la
amplitud de las dos reconstrucciones, en el 80\,\% central de la placa: la
medida de la interferometría holográfica.

\begin{table}[htbp]
  \centering
  \small
  \caption[Interferometría simulada: error de la ida]{Error relativo del
  campo en el sensor al simular el holograma, frente a la referencia, por
  cociente $q$ y método. Con 2048 píxeles, cada $q$ cae al doble de $z$. Los
  valores menores que $10^{-8}$ están en el piso de la referencia. Abajo, el
  mayor arreglo de cada método y su tiempo por propagación a $q = 20$, en CPU
  y en complex128. El MPASM con $s = 4$ da lo mismo que con $s = 2$ en todos
  los $q$, con \SI{268}{MB} y unos \SI{3.4}{s}, y no se muestra.}
  \label{tab:larga}
  \setlength{\tabcolsep}{4pt}
  \begin{tabular}{rrcccccccc}
    \toprule
    & & & \multicolumn{3}{c}{FFT-ASM, relleno} & BLAS
      & \multicolumn{2}{c}{MPASM} \\
    \cmidrule(lr){4-6}\cmidrule(lr){8-9}
    $q$ & $z$ [mm] & $K_f$ ($s = 1$) & $\times1$ & $\times2$ & $\times4$
      & $\times2$ & $s = 1$ & $s = 2$ \\
    \midrule
    0.5 & 11.5  & 1.00 & $5.2\times10^{-4}$ & $<10^{-8}$ & $<10^{-8}$ & $<10^{-8}$
        & $5.2\times10^{-4}$ & $<10^{-8}$ \\
    1   & 22.9  & 1.00 & $1.8\times10^{-2}$ & $<10^{-8}$ & $<10^{-8}$ & $<10^{-8}$
        & $1.6\times10^{-2}$ & $<10^{-8}$ \\
    2   & 45.8  & 1.42 & 0.11 & $<10^{-8}$ & $<10^{-8}$ & $<10^{-8}$ & $<10^{-8}$
        & $<10^{-8}$ \\
    3   & 68.7  & 1.73 & 0.20 & $<10^{-8}$ & $<10^{-8}$ & $<10^{-8}$ & $<10^{-8}$
        & $<10^{-8}$ \\
    5   & 114.6 & 2.24 & 0.38 & $<10^{-8}$ & $<10^{-8}$ & $<10^{-8}$ & $<10^{-8}$
        & $<10^{-8}$ \\
    10  & 229.1 & 3.16 & 0.82 & $1.8\times10^{-5}$ & $<10^{-8}$ & $3.6\times10^{-5}$
        & $<10^{-8}$ & $<10^{-8}$ \\
    20  & 458.2 & 4.47 & 2.05 & 0.24 & $<10^{-8}$ & 0.094 & $7.4\times10^{-7}$
        & $<10^{-8}$ \\
    \midrule
    \multicolumn{3}{l}{mayor arreglo [MB]} & 17 & 67 & 268 & 67 & 17 & 67 \\
    \multicolumn{3}{l}{tiempo a $q = 20$ [s]} & 0.17 & 0.73 & 3.4 & 0.79 & 0.37
      & 1.07 \\
    \bottomrule
  \end{tabular}
\end{table}

\paragraph{Al propagar.} La Tabla~\ref{tab:larga} y la
Figura~\ref{fig:larga}a dan el error de la ida. El $K_f$ automático vale 1
hasta $q = 1$ y desde ahí crece como $\sqrt{q}$, hasta 4.47 en $q = 20$. El
FFT-ASM necesita un relleno que crece con la distancia. Sin relleno ya se
aparta de la referencia a $q = 0.5$, y a $q = 20$ su error es de 2, el doble de
la señal. Con relleno al doble es exacto hasta $q = 5$, se aparta del piso a
$q = 10$ y a $q = 20$ da 0.24. Solo con relleno al cuádruple llega exacto a
$q = 20$, con un arreglo de \SI{268}{MB} y \SI{3.4}{s} por propagación. El
MPASM con $s = 1$, que no tiene relleno, se comporta como el FFT-ASM sin
relleno mientras $K_f = 1$ ---a $q \le 1$ necesita $s = 2$---, pero desde
$q = 2$, cuando $K_f$ pasa de 1, queda en el piso de la referencia, y a
$q = 20$ da $7\times10^{-7}$ con \SI{17}{MB} y \SI{0.37}{s}. El BLAS no se
contamina de \emph{aliasing}, pero a $q = 20$ recorta la banda y se equivoca
en 0.094. Es la afirmación de la discusión, medida: el $K_f$ que crece con la
distancia hace el papel del relleno que el FFT-ASM necesita, y lo hace sin
agrandar la malla.

\paragraph{Al medir la deformación.} La Figura~\ref{fig:larga}b da el RMS de
la diferencia de fase reconstruida, y aquí el propagador no decide. La
reconstrucción ideal ---la referencia devuelta desde el holograma--- se aparta
de la deformación 0.009~rad a $q = 2$ y 0.33~rad a $q = 20$, y ninguno de los
métodos sobre la malla del objeto se aparta de ella más de 0.003~rad, ni
siquiera el FFT-ASM sin relleno. Lo que limita la medida es la apertura del
sensor: a $q = 20$ solo le llega el 19\,\% de la luz del objeto, y la
reconstrucción resuelve unos \SI{69}{\micro\metre}, $\lambda z/(N\delta)$, una
escala en la que el gradiente de la deformación ya pesa
(Figura~\ref{fig:larga}c y~d). La transformada de Fresnel da la misma
diferencia de fase desde $q = 1$, con el menor costo (\SI{0.17}{s}), y falla a
$q = 0.5$ (1.27~rad), donde su \emph{chirp} de entrada queda submuestreado.

\begin{figure}[htbp]
  \centering
  \includegraphics[width=\textwidth]{interferometria.pdf}
  \caption[Interferometría holográfica simulada a distancias
  largas]{Interferometría holográfica simulada: un difusor de \SI{3.2}{mm}
  ante un sensor de $1024\times1024$ píxeles de \SI{3.45}{\micro\metre}, a
  \SI{532}{nm}. (a)~Error del campo en el sensor al simular el holograma,
  contra el cociente $q$; la banda gris es el piso de la referencia.
  (b)~RMS de la diferencia de fase reconstruida frente a la deformación: las
  curvas de los métodos quedan bajo la de la reconstrucción ideal. (c)~La
  diferencia de fase de la deformación y (d)~la reconstruida con el MPASM,
  $s = 1$, a $q = 20$ (\SI{0.46}{m}).}
  \label{fig:larga}
\end{figure}

\paragraph{Lo que dice el ensayo.} En montajes de propagación larga, el MPASM
rinde para propagar: simula el holograma con la exactitud del FFT-ASM con
cuatro veces el relleno y con dieciséis veces menos memoria, y lo hace porque
su $K_f$ crece con la distancia. Para la medida de la interferometría
holográfica, en cambio, el propagador no decide: la fija la apertura del
sensor, y la transformada de Fresnel, que es la de uso
corriente~\cite{schnars2005}, basta. Es la misma separación entre propagar y
reconstruir del resto del trabajo. El ensayo tiene los límites de una
simulación: un solo objeto; un difusor de banda limitada ---uno real se abre
más, lo que solo adelanta el desborde de la ventana---; el campo complejo en el
sensor, como lo entrega el desplazamiento de fase, sin onda de referencia
fuera de eje ni ruido, y el cálculo en CPU.
